\chapter{Chapter07 — Quadratic Circuit Reduction}
%%% generated by the blueprint pipeline from FABL.Chapter07.QuadraticCircuitReduction
\section{Overview}

Declarations of the Lean module \texttt{FABL.Chapter07.QuadraticCircuitReduction}.

\section{Declarations}

\begin{definition}[Quadratic equation from indexed terms]
\lean{FABL.homogeneousQuadraticEquationOfTerms}
\statementsource{Analysis-of-Boolean-Functions-by-Ryan-ODonnell}{pdf-ffbe9c7b3306-p198-b003}
\label{FABL.homogeneousQuadraticEquationOfTerms}
\leanok
\uses{FABL.HomogeneousQuadraticEquation, FABL.𝔽₂}
Let $N$ be a natural number, let $T$ be a finite list of ordered pairs $(i,j)$ with
$i,j\in\{0,\ldots,N-1\}$, and let $b\in\mathbb{F}_2$. The associated homogeneous
quadratic equation over $\mathbb{F}_2$ in $N$ variables has right-hand side $b$ and
coefficient matrix obtained by summing, with multiplicity, the matrices that have a $1$
in position $(i,j)$ and zeros elsewhere for each pair $(i,j)$ in $T$.
\end{definition}

\begin{theorem}[Pairing with a Single-Entry Matrix]
\lean{FABL.binaryMatrixPairing_single}
\label{FABL.binaryMatrixPairing_single}
\leanok
\uses{FABL.F₂Cube, FABL.binaryMatrixPairing, FABL.rankOneSelfMatrix}
\difficulty{4}
Let $N$ be a natural number, let $i,j$ be indices in $\{0,\ldots,N-1\}$, and let $z$ be
a vector of length $N$ over $\mathbb{F}_2$. Pairing the matrix with entry $1$ at $(i,j)$
and $0$ elsewhere with the matrix whose $(a,b)$ entry is $z_a z_b$ gives $z_i z_j$.
\end{theorem}

\begin{theorem}[Additivity in the First Matrix Argument]
\lean{FABL.binaryMatrixPairing_add_left}
\label{FABL.binaryMatrixPairing_add_left}
\leanok
\uses{FABL.F₂Cube, FABL.binaryMatrixPairing, FABL.rankOneSelfMatrix, FABL.𝔽₂}
\difficulty{1}
Let $N$ be a natural number, let $A$ and $B$ be $N\times N$ matrices over
$\mathbb{F}_2$, and let $z$ be a vector of length $N$ over $\mathbb{F}_2$. Pairing $A+B$
with the matrix whose $(i,j)$ entry is $z_i z_j$ gives the sum of the pairings of $A$
and $B$ with that matrix.
\end{theorem}

\begin{theorem}[Pairing a Sum of Single-Entry Matrices]
\lean{FABL.binaryMatrixPairing_sum_single}
\label{FABL.binaryMatrixPairing_sum_single}
\leanok
\uses{FABL.F₂Cube, FABL.binaryMatrixPairing, FABL.binaryMatrixPairing_add_left, FABL.binaryMatrixPairing_single, FABL.rankOneSelfMatrix}
\difficulty{4}
Let $N$ be a natural number, let $\mathrm{terms}$ be a finite list of pairs of indices
in $\{0,\ldots,N-1\}$, and let $z$ be a vector of length $N$ over $\mathbb{F}_2$.
Pairing the sum of the single-entry matrices specified by the listed pairs with the
matrix whose $(i,j)$ entry is $z_i z_j$ equals the sum, over the list, of $z_i z_j$ for
each listed pair $(i,j)$.
\end{theorem}

\begin{theorem}[Quadratic equation from a list of terms]
\lean{FABL.homogeneousQuadraticEquationOfTerms_satisfiedBy_iff}
\label{FABL.homogeneousQuadraticEquationOfTerms_satisfiedBy_iff}
\leanok
\uses{FABL.F₂Cube, FABL.HomogeneousQuadraticEquation.SatisfiedBy, FABL.binaryMatrixPairing_sum_single, FABL.homogeneousQuadraticEquationOfTerms, FABL.𝔽₂}
\difficulty{3}
Let $N$ be a natural number, let $\mathbb F_2$ be the field with two elements, take a
finite list of ordered pairs of coordinates in $\{0,\ldots,N-1\}$, choose a right-hand
side $c\in\mathbb F_2$, and let $z\in\mathbb F_2^N$. The quadratic equation formed from
that list is satisfied by $z$ if and only if the sum over the list, counting
repetitions, of the products $z_a z_b$ for its pairs $(a,b)$ equals $c$ in $\mathbb
F_2$.
\end{theorem}

\begin{definition}[Circuit quadratic variable count]
\lean{FABL.BooleanCircuit.circuitQuadraticVariableCount}
\statementsource{Analysis-of-Boolean-Functions-by-Ryan-ODonnell}{pdf-ffbe9c7b3306-p198-b004}
\label{FABL.BooleanCircuit.circuitQuadraticVariableCount}
\leanok
\uses{FABL.BooleanCircuit}
For a Boolean circuit with $n$ input bits and $g$ gates, define its quadratic variable
count to be $1+n+g$.
\end{definition}

\begin{definition}[Constant variable index]
\lean{FABL.BooleanCircuit.circuitQuadraticConstantIndex}
\label{FABL.BooleanCircuit.circuitQuadraticConstantIndex}
\leanok
\uses{FABL.BooleanCircuit, FABL.BooleanCircuit.circuitQuadraticVariableCount}
For a Boolean circuit $C$ with $n$ input bits and $m$ gates, index the variables by the
disjoint union of one constant variable and $n+m$ further variables. The constant
variable has index $0$ in this ordering.
\end{definition}

\begin{definition}[Input variable index]
\lean{FABL.BooleanCircuit.circuitQuadraticInputIndex}
\label{FABL.BooleanCircuit.circuitQuadraticInputIndex}
\leanok
\uses{FABL.BooleanCircuit, FABL.BooleanCircuit.circuitQuadraticVariableCount}
For a Boolean circuit $C$ with $n$ input bits and $m$ gates, index the variables by the
disjoint union of one constant variable, the $n$ input variables, and the $m$ gate
variables. The variable corresponding to input bit $i$ has index $i+1$.
\end{definition}

\begin{definition}[Quadratic gate coordinate]
\lean{FABL.BooleanCircuit.circuitQuadraticGateIndex}
\statementsource{Analysis-of-Boolean-Functions-by-Ryan-ODonnell}{pdf-ffbe9c7b3306-p198-b004}
\label{FABL.BooleanCircuit.circuitQuadraticGateIndex}
\leanok
\uses{FABL.BooleanCircuit, FABL.BooleanCircuit.circuitQuadraticVariableCount}
Let $C$ be a Boolean circuit with $n$ input bits and $g$ gates, and let $i$ be a gate
index. Among the $1+n+g$ coordinates consisting of one constant coordinate, the $n$
input coordinates, and the $g$ gate coordinates, the coordinate assigned to gate $i$ is
the $(1+n+i)$-th coordinate, with indexing starting at zero.
\end{definition}

\begin{definition}[Quadratic circuit witness]
\lean{FABL.BooleanCircuit.circuitQuadraticWitness}
\statementsource{Analysis-of-Boolean-Functions-by-Ryan-ODonnell}{pdf-ffbe9c7b3306-p198-b004}
\label{FABL.BooleanCircuit.circuitQuadraticWitness}
\leanok
\uses{FABL.BooleanCircuit, FABL.BooleanCircuit.circuitQuadraticVariableCount, FABL.BooleanCircuit.evalGate, FABL.F₂Cube, FABL.pcppBoolF₂Equiv}
Let $C$ be a Boolean circuit with $n$ input bits and $g$ gates, and let $w$ be an
assignment of Boolean values to its input bits. Form a vector over $\mathbb{F}_2$ with
$1+n+g$ coordinates: its constant coordinate is $1$, its input coordinates encode the
values of $w$, and its gate coordinates encode the values produced by evaluating the
corresponding gates of $C$ on $w$.
\end{definition}

\begin{theorem}[Constant coordinate of the witness]
\lean{FABL.BooleanCircuit.circuitQuadraticWitness_constant}
\label{FABL.BooleanCircuit.circuitQuadraticWitness_constant}
\leanok
\uses{FABL.BooleanCircuit, FABL.BooleanCircuit.circuitQuadraticConstantIndex, FABL.BooleanCircuit.circuitQuadraticWitness}
\difficulty{1}
For every Boolean circuit $C$ with $n$ input bits and every Boolean input $w$, the
quadratic witness of $C$ at $w$ has value $1$ in the field with two elements at its
constant coordinate.
\end{theorem}

\begin{theorem}[Witness at an input coordinate]
\lean{FABL.BooleanCircuit.circuitQuadraticWitness_input}
\label{FABL.BooleanCircuit.circuitQuadraticWitness_input}
\leanok
\uses{FABL.BooleanCircuit, FABL.BooleanCircuit.circuitQuadraticInputIndex, FABL.BooleanCircuit.circuitQuadraticWitness, FABL.pcppBoolF₂Equiv}
\difficulty{1}
Let $C$ be a Boolean circuit on $n$ input bits, let $w$ be an input assignment, and let
$i$ be an input position. The circuit witness associated with $w$, evaluated at the
coordinate for input $i$, equals the element of $mathbb{F}_2$ corresponding to the bit
$w(i)$.
\end{theorem}

\begin{theorem}[Gate coordinate of the witness]
\lean{FABL.BooleanCircuit.circuitQuadraticWitness_gate}
\label{FABL.BooleanCircuit.circuitQuadraticWitness_gate}
\leanok
\uses{FABL.BooleanCircuit, FABL.BooleanCircuit.circuitQuadraticGateIndex, FABL.BooleanCircuit.circuitQuadraticWitness, FABL.BooleanCircuit.evalGate, FABL.pcppBoolF₂Equiv}
\difficulty{1}
For every Boolean circuit $C$ with $n$ input bits, every Boolean input $w$, and every
gate $i$ of $C$, the quadratic witness of $C$ at $w$, evaluated at the coordinate
assigned to gate $i$, equals the element of the field with two elements corresponding to
the Boolean value of gate $i$ on input $w$.
\end{theorem}

\begin{theorem}[Field encoding of Boolean disjunction]
\lean{FABL.BooleanCircuit.pcppBoolF₂Equiv_or}
\leanok
\uses{FABL.pcppBoolF₂Equiv}
\difficulty{4}
For Boolean values $a$ and $b$, the image of their disjunction in $\mathbb{F}_2$ equals
$\bar a+\bar b+\bar a\bar b$, where $\bar a$ and $\bar b$ are their images in
$\mathbb{F}_2$.
\end{theorem}

\begin{definition}[Quadratic constant equation]
\lean{FABL.BooleanCircuit.circuitQuadraticConstantEquation}
\statementsource{Analysis-of-Boolean-Functions-by-Ryan-ODonnell}{pdf-ffbe9c7b3306-p198-b004}
\label{FABL.BooleanCircuit.circuitQuadraticConstantEquation}
\leanok
\uses{FABL.BooleanCircuit, FABL.BooleanCircuit.circuitQuadraticConstantIndex, FABL.BooleanCircuit.circuitQuadraticVariableCount, FABL.HomogeneousQuadraticEquation, FABL.homogeneousQuadraticEquationOfTerms}
Let $C$ be a Boolean circuit with $n$ input bits and $g$ gates, and use $1+n+g$
variables over $\mathbb{F}_2$. The associated homogeneous quadratic equation has the
single quadratic term given by the square of the constant-coordinate variable and has
right-hand side $1$.
\end{definition}

\begin{definition}[Quadratic gate equation]
\lean{FABL.BooleanCircuit.circuitQuadraticGateEquation}
\statementsource{Analysis-of-Boolean-Functions-by-Ryan-ODonnell}{pdf-ffbe9c7b3306-p198-b004}
\label{FABL.BooleanCircuit.circuitQuadraticGateEquation}
\leanok
\uses{FABL.BooleanCircuit, FABL.BooleanCircuit.circuitQuadraticConstantIndex, FABL.BooleanCircuit.circuitQuadraticGateIndex, FABL.BooleanCircuit.circuitQuadraticInputIndex, FABL.BooleanCircuit.circuitQuadraticVariableCount, FABL.BooleanCircuit.previousGateIndex, FABL.HomogeneousQuadraticEquation, FABL.homogeneousQuadraticEquationOfTerms}
Let $C$ be a Boolean circuit with $n$ input bits and $g$ gates, and use $1+n+g$
variables over $\mathbb{F}_2$. For a gate $i$, define a homogeneous quadratic equation
with right-hand side $0$ as follows, writing $c$ for the constant-coordinate variable,
$x_j$ for input variables, and $y_k$ for gate variables: if gate $i$ reads input $j$,
its equation is $y_i c+x_j c=0$; if it is the constant $0$, its equation is $y_i c=0$;
if it is the constant $1$, its equation is $y_i c+c^2=0$; if it negates an earlier gate
$k$, its equation is $y_i c+y_k c+c^2=0$; if it is the conjunction of earlier gates $k$
and $$, its equation is $y_i c+y_k y_\ell=0$; and if it is their disjunction, its
equation is $y_i c+y_k c+y_\ell c+y_k y_\ell=0$.
\end{definition}

\begin{definition}[Quadratic output equation]
\lean{FABL.BooleanCircuit.circuitQuadraticOutputEquation}
\statementsource{Analysis-of-Boolean-Functions-by-Ryan-ODonnell}{pdf-ffbe9c7b3306-p198-b004}
\label{FABL.BooleanCircuit.circuitQuadraticOutputEquation}
\leanok
\uses{FABL.BooleanCircuit, FABL.BooleanCircuit.circuitQuadraticConstantIndex, FABL.BooleanCircuit.circuitQuadraticGateIndex, FABL.BooleanCircuit.circuitQuadraticVariableCount, FABL.HomogeneousQuadraticEquation, FABL.homogeneousQuadraticEquationOfTerms}
Let $C$ be a Boolean circuit with $n$ input bits and $g$ gates, and use $1+n+g$
variables over $\mathbb{F}_2$. The associated homogeneous quadratic equation has
right-hand side $0$ and is $y_o c+c^2=0$, where $y_o$ is the variable for the
distinguished output gate and $c$ is the constant-coordinate variable.
\end{definition}

\begin{definition}[Quadratic circuit equation system]
\lean{FABL.BooleanCircuit.circuitQuadraticSystem}
\statementsource{Analysis-of-Boolean-Functions-by-Ryan-ODonnell}{pdf-ffbe9c7b3306-p198-b004}
\label{FABL.BooleanCircuit.circuitQuadraticSystem}
\leanok
\uses{FABL.BooleanCircuit, FABL.BooleanCircuit.circuitQuadraticConstantEquation, FABL.BooleanCircuit.circuitQuadraticGateEquation, FABL.BooleanCircuit.circuitQuadraticOutputEquation, FABL.BooleanCircuit.circuitQuadraticVariableCount, FABL.HomogeneousQuadraticSystem}
Let $C$ be a Boolean circuit with $n$ input bits and $g$ gates, and use $1+n+g$
variables over $\mathbb{F}_2$. The associated system, in order, consists of the
constant equation, one gate equation for each gate in the circuit's order, and the
output equation.
\end{definition}

\begin{definition}[Input recovered from a quadratic witness]
\lean{FABL.BooleanCircuit.circuitQuadraticInput}
\statementsource{Analysis-of-Boolean-Functions-by-Ryan-ODonnell}{pdf-ffbe9c7b3306-p198-b004}
\label{FABL.BooleanCircuit.circuitQuadraticInput}
\leanok
\uses{FABL.BooleanCircuit, FABL.BooleanCircuit.circuitQuadraticInputIndex, FABL.BooleanCircuit.circuitQuadraticVariableCount, FABL.F₂Cube, FABL.pcppBoolF₂Equiv}
Let $C$ be a Boolean circuit with $n$ input bits and $g$ gates, and let $z$ be a vector
over $\mathbb{F}_2$ with $1+n+g$ coordinates. The Boolean input assignment recovered
from $z$ assigns to each input position the Boolean value encoded by the corresponding
input coordinate of $z$.
\end{definition}

\begin{theorem}[Boolean input encoding identity]
\lean{FABL.BooleanCircuit.pcppBoolF₂Equiv_circuitQuadraticInput}
\leanok
\uses{FABL.BooleanCircuit, FABL.BooleanCircuit.circuitQuadraticInput, FABL.BooleanCircuit.circuitQuadraticInputIndex, FABL.BooleanCircuit.circuitQuadraticVariableCount, FABL.F₂Cube, FABL.pcppBoolF₂Equiv}
\difficulty{1}
Let $C$ be a Boolean circuit with $n$ input bits and $g$ gates, and let $z$ assign an
element of $[?]$ to each of its $1+n+g$ coordinates. Identify $\text{false}$ with $0$
and $\text{true}$ with $1$ in the field $[?]$. For every input index $i$, encoding the
Boolean value obtained by decoding the coordinate of $z$ assigned to input $i$ recovers
that coordinate.
\end{theorem}

\begin{theorem}[Constant equation holds for the circuit witness]
\lean{FABL.BooleanCircuit.circuitQuadraticConstantEquation_witness}
\label{FABL.BooleanCircuit.circuitQuadraticConstantEquation_witness}
\leanok
\uses{FABL.BooleanCircuit, FABL.BooleanCircuit.circuitQuadraticConstantEquation, FABL.BooleanCircuit.circuitQuadraticWitness, FABL.HomogeneousQuadraticEquation.SatisfiedBy, FABL.homogeneousQuadraticEquationOfTerms_satisfiedBy_iff}
\difficulty{3}
Let $C$ be a Boolean circuit on $n$ input bits, and let $w$ be an input assignment. Form
the binary vector whose constant coordinate is $1$, whose input coordinates encode $w$,
and whose gate coordinates encode the values of the gates on $w$. This vector satisfies
the homogeneous quadratic equation that fixes the constant coordinate to $1$.
\end{theorem}

\begin{theorem}[Circuit gate equations hold on the witness]
\lean{FABL.BooleanCircuit.circuitQuadraticGateEquation_witness}
\label{FABL.BooleanCircuit.circuitQuadraticGateEquation_witness}
\leanok
\uses{FABL.BooleanCircuit, FABL.BooleanCircuit.circuitQuadraticGateEquation, FABL.BooleanCircuit.circuitQuadraticWitness, FABL.BooleanCircuit.circuitQuadraticWitness_constant, FABL.BooleanCircuit.circuitQuadraticWitness_gate, FABL.BooleanCircuit.circuitQuadraticWitness_input, FABL.BooleanCircuit.evalGate, FABL.BooleanCircuit.evalGate_eq_and, FABL.BooleanCircuit.evalGate_eq_const, FABL.BooleanCircuit.evalGate_eq_input, FABL.BooleanCircuit.evalGate_eq_neg, FABL.BooleanCircuit.evalGate_eq_or, FABL.BooleanCircuit.pcppBoolF₂Equiv_or, FABL.BooleanCircuit.previousGateIndex, FABL.HomogeneousQuadraticEquation.SatisfiedBy, FABL.homogeneousQuadraticEquationOfTerms_satisfiedBy_iff, FABL.pcppBoolF₂Equiv_and, FABL.pcppBoolF₂Equiv_not}
\difficulty{5}
Let $C$ be a Boolean circuit on $n$ input bits, let $w$ be a Boolean input, and let $i$
be a gate index. The homogeneous quadratic equation encoding gate $i$ is satisfied by
the quadratic witness obtained from $w$ and the evaluations of the circuit’s gates.
\end{theorem}

\begin{theorem}[Output equation for an accepting witness]
\lean{FABL.BooleanCircuit.circuitQuadraticOutputEquation_witness}
\label{FABL.BooleanCircuit.circuitQuadraticOutputEquation_witness}
\leanok
\uses{FABL.BooleanCircuit, FABL.BooleanCircuit.acceptedProperty, FABL.BooleanCircuit.circuitQuadraticOutputEquation, FABL.BooleanCircuit.circuitQuadraticWitness, FABL.BooleanCircuit.circuitQuadraticWitness_constant, FABL.BooleanCircuit.circuitQuadraticWitness_gate, FABL.BooleanCircuit.eval, FABL.HomogeneousQuadraticEquation.SatisfiedBy, FABL.homogeneousQuadraticEquationOfTerms_satisfiedBy_iff}
\difficulty{3}
Let $C$ be a Boolean circuit on $n$ input bits, and let $w$ be an input word that the
circuit accepts. The canonical quadratic witness formed from $w$ and the values of all
gates satisfies the circuit’s quadratic output equation.
\end{theorem}

\begin{theorem}[Circuit witness satisfies the quadratic system]
\lean{FABL.BooleanCircuit.circuitQuadraticWitness_satisfies}
\proofsource{Analysis-of-Boolean-Functions-by-Ryan-ODonnell}{pdf-ffbe9c7b3306-p198-b004}
\label{FABL.BooleanCircuit.circuitQuadraticWitness_satisfies}
\leanok
\uses{FABL.BooleanCircuit, FABL.BooleanCircuit.acceptedProperty, FABL.BooleanCircuit.circuitQuadraticConstantEquation_witness, FABL.BooleanCircuit.circuitQuadraticGateEquation_witness, FABL.BooleanCircuit.circuitQuadraticOutputEquation_witness, FABL.BooleanCircuit.circuitQuadraticSystem, FABL.BooleanCircuit.circuitQuadraticWitness, FABL.SatisfiesHomogeneousQuadraticSystem}
\difficulty{4}
Let $C$ be a Boolean circuit on $n$ input bits and let $w$ be an input assignment
accepted by $C$. Then the circuit witness associated with $w$ satisfies the homogeneous
quadratic system encoding $C$.
\end{theorem}

\begin{theorem}[Constant coordinate of a satisfying assignment]
\lean{FABL.BooleanCircuit.circuitQuadraticConstant_eq_one_of_satisfies}
\label{FABL.BooleanCircuit.circuitQuadraticConstant_eq_one_of_satisfies}
\leanok
\uses{FABL.BooleanCircuit, FABL.BooleanCircuit.circuitQuadraticConstantEquation, FABL.BooleanCircuit.circuitQuadraticConstantIndex, FABL.BooleanCircuit.circuitQuadraticSystem, FABL.BooleanCircuit.circuitQuadraticVariableCount, FABL.F₂Cube, FABL.SatisfiesHomogeneousQuadraticSystem, FABL.homogeneousQuadraticEquationOfTerms_satisfiedBy_iff}
\difficulty{3}
Let $C$ be a Boolean circuit on $n$ input bits, and let $z$ be an assignment over $$ to
the circuit’s quadratic variables. If $z$ satisfies every equation in the circuit’s
associated homogeneous quadratic system, then its constant coordinate is $1$.
\end{theorem}

\begin{theorem}[Gate values in satisfying assignments]
\lean{FABL.BooleanCircuit.circuitQuadraticGateValue_eq_evalGate_of_satisfies}
\proofsource{Analysis-of-Boolean-Functions-by-Ryan-ODonnell}{pdf-ffbe9c7b3306-p198-b004}
\label{FABL.BooleanCircuit.circuitQuadraticGateValue_eq_evalGate_of_satisfies}
\leanok
\uses{FABL.BooleanCircuit, FABL.BooleanCircuit.circuitQuadraticConstant_eq_one_of_satisfies, FABL.BooleanCircuit.circuitQuadraticGateEquation, FABL.BooleanCircuit.circuitQuadraticGateIndex, FABL.BooleanCircuit.circuitQuadraticInput, FABL.BooleanCircuit.circuitQuadraticInputIndex, FABL.BooleanCircuit.circuitQuadraticSystem, FABL.BooleanCircuit.circuitQuadraticVariableCount, FABL.BooleanCircuit.evalGate, FABL.BooleanCircuit.evalGate_eq_and, FABL.BooleanCircuit.evalGate_eq_const, FABL.BooleanCircuit.evalGate_eq_input, FABL.BooleanCircuit.evalGate_eq_neg, FABL.BooleanCircuit.evalGate_eq_or, FABL.BooleanCircuit.pcppBoolF₂Equiv_circuitQuadraticInput, FABL.BooleanCircuit.pcppBoolF₂Equiv_or, FABL.BooleanCircuit.previousGateIndex, FABL.F₂Cube, FABL.SatisfiesHomogeneousQuadraticSystem, FABL.homogeneousQuadraticEquationOfTerms_satisfiedBy_iff, FABL.pcppBoolF₂Equiv, FABL.pcppBoolF₂Equiv_and, FABL.pcppBoolF₂Equiv_not}
\difficulty{6}
Let $C$ be a Boolean circuit on $n$ input bits, and let $z$ assign an element of
$\mathbb{F}_2$ to each of its $1+n+g$ variables, where $g$ is the number of gates. If
$z$ satisfies the circuit’s homogeneous quadratic system, then for every gate $i$, its
variable equals the $\mathbb{F}_2$ encoding of the Boolean value of gate $i$ on the
input decoded from the input variables of $z$.
\end{theorem}

\begin{theorem}[Acceptance from a satisfying quadratic assignment]
\lean{FABL.BooleanCircuit.circuitQuadraticInput_accepted_of_satisfies}
\label{FABL.BooleanCircuit.circuitQuadraticInput_accepted_of_satisfies}
\leanok
\uses{FABL.BooleanCircuit, FABL.BooleanCircuit.acceptedProperty, FABL.BooleanCircuit.circuitQuadraticConstant_eq_one_of_satisfies, FABL.BooleanCircuit.circuitQuadraticGateValue_eq_evalGate_of_satisfies, FABL.BooleanCircuit.circuitQuadraticInput, FABL.BooleanCircuit.circuitQuadraticOutputEquation, FABL.BooleanCircuit.circuitQuadraticSystem, FABL.BooleanCircuit.circuitQuadraticVariableCount, FABL.BooleanCircuit.eval, FABL.F₂Cube, FABL.SatisfiesHomogeneousQuadraticSystem, FABL.homogeneousQuadraticEquationOfTerms_satisfiedBy_iff, FABL.pcppBoolF₂Equiv}
\difficulty{4}
Let $C$ be a Boolean circuit on $n$ input bits, and let $z$ assign an element of
$\mathbb{F}_2$ to each of its $1+n+g$ variables, where $g$ is the number of gates. If
$z$ satisfies the circuit’s homogeneous quadratic system, then the circuit accepts the
input obtained by decoding the input variables of $z$.
\end{theorem}

\begin{theorem}[Circuit acceptance and quadratic extensions]
\lean{FABL.BooleanCircuit.acceptedProperty_iff_exists_satisfyingQuadraticExtension}
\proofsource{Analysis-of-Boolean-Functions-by-Ryan-ODonnell}{pdf-ffbe9c7b3306-p198-b004}
\label{FABL.BooleanCircuit.acceptedProperty_iff_exists_satisfyingQuadraticExtension}
\leanok
\uses{FABL.BooleanCircuit, FABL.BooleanCircuit.acceptedProperty, FABL.BooleanCircuit.circuitQuadraticInput, FABL.BooleanCircuit.circuitQuadraticInputIndex, FABL.BooleanCircuit.circuitQuadraticInput_accepted_of_satisfies, FABL.BooleanCircuit.circuitQuadraticSystem, FABL.BooleanCircuit.circuitQuadraticVariableCount, FABL.BooleanCircuit.circuitQuadraticWitness, FABL.BooleanCircuit.circuitQuadraticWitness_input, FABL.BooleanCircuit.circuitQuadraticWitness_satisfies, FABL.BooleanCircuit.pcppBoolF₂Equiv_circuitQuadraticInput, FABL.F₂Cube, FABL.SatisfiesHomogeneousQuadraticSystem, FABL.pcppBoolF₂Equiv}
\difficulty{4}
Let $C$ be a Boolean circuit on $n$ input bits, and let $w$ be an input assignment. The
circuit accepts $w$ if and only if there is a vector over $\bbr_2$ with one coordinate
for the constant and one for each input and gate that satisfies the circuit's
homogeneous quadratic system and whose input coordinates encode $w$.
\end{theorem}

\begin{theorem}[Circuit satisfiability and quadratic systems]
\lean{FABL.BooleanCircuit.satisfiable_iff_exists_satisfiesCircuitQuadraticSystem}
\proofsource{Analysis-of-Boolean-Functions-by-Ryan-ODonnell}{pdf-ffbe9c7b3306-p198-b004}
\label{FABL.BooleanCircuit.satisfiable_iff_exists_satisfiesCircuitQuadraticSystem}
\leanok
\uses{FABL.BooleanCircuit, FABL.BooleanCircuit.Satisfiable, FABL.BooleanCircuit.circuitQuadraticInput, FABL.BooleanCircuit.circuitQuadraticInput_accepted_of_satisfies, FABL.BooleanCircuit.circuitQuadraticSystem, FABL.BooleanCircuit.circuitQuadraticVariableCount, FABL.BooleanCircuit.circuitQuadraticWitness, FABL.BooleanCircuit.circuitQuadraticWitness_satisfies, FABL.F₂Cube, FABL.SatisfiesHomogeneousQuadraticSystem}
\difficulty{1}
Let $C$ be a Boolean circuit with $n$ input bits and gates ordered so that each gate
uses only input bits or earlier gates. Form a homogeneous quadratic system over the
field with two elements, with one coordinate for a constant, each input bit, and each
gate output; impose that the constant coordinate is $1$, that every gate coordinate
agrees with its gate operation, and that the output coordinate is $1$. Then some input
makes $C$ output true if and only if this system has a solution.
\end{theorem}

\begin{theorem}[Length of the quadratic system]
\lean{FABL.BooleanCircuit.circuitQuadraticSystem_length}
\label{FABL.BooleanCircuit.circuitQuadraticSystem_length}
\leanok
\uses{FABL.BooleanCircuit, FABL.BooleanCircuit.circuitQuadraticSystem}
\difficulty{1}
For every Boolean circuit $C$ with $n$ input bits and $g$ gates, its associated
homogeneous quadratic system has exactly $g+2$ equations.
\end{theorem}

\begin{theorem}[Quadratic variable count formula]
\lean{FABL.BooleanCircuit.circuitQuadraticVariableCount_eq}
\label{FABL.BooleanCircuit.circuitQuadraticVariableCount_eq}
\leanok
\uses{FABL.BooleanCircuit, FABL.BooleanCircuit.circuitQuadraticVariableCount}
\difficulty{1}
For every Boolean circuit $C$ with $n$ input bits and $g$ gates, its quadratic variable
count is $n+g+1$.
\end{theorem}

\begin{definition}[Dense system size]
\lean{FABL.BooleanCircuit.circuitQuadraticSystemDenseSize}
\label{FABL.BooleanCircuit.circuitQuadraticSystemDenseSize}
\leanok
\uses{FABL.BooleanCircuit, FABL.BooleanCircuit.circuitQuadraticSystem, FABL.BooleanCircuit.circuitQuadraticVariableCount}
For a Boolean circuit $C$, let $L$ be the length of its associated quadratic system and
let $V$ be its quadratic variable count. The dense size of the system is $L V^2$.
\end{definition}

\begin{theorem}[Dense size of the quadratic system]
\lean{FABL.BooleanCircuit.circuitQuadraticSystemDenseSize_le_cube}
\label{FABL.BooleanCircuit.circuitQuadraticSystemDenseSize_le_cube}
\leanok
\uses{FABL.BooleanCircuit, FABL.BooleanCircuit.circuitQuadraticSystemDenseSize, FABL.BooleanCircuit.circuitQuadraticSystem_length, FABL.BooleanCircuit.circuitQuadraticVariableCount, FABL.BooleanCircuit.circuitQuadraticVariableCount_eq}
\difficulty{4}
For every Boolean circuit $C$ with $n$ input bits and $g$ gates, the dense size of its
associated homogeneous quadratic system is at most $(n+g+2)^3$.
\end{theorem}

\begin{theorem}[Satisfiability with no input bits]
\lean{FABL.BooleanCircuit.zeroInput_satisfiable_iff}
\label{FABL.BooleanCircuit.zeroInput_satisfiable_iff}
\leanok
\uses{FABL.BooleanCircuit, FABL.BooleanCircuit.Satisfiable, FABL.BooleanCircuit.acceptedProperty}
\difficulty{4}
Let $C$ be a Boolean circuit with no input bits. It is satisfiable if and only if it
accepts the unique assignment on the empty set of input coordinates.
\end{theorem}

\begin{definition}[Quadratic circuit proof length]
\lean{FABL.BooleanCircuit.circuitQuadraticPCPPProofLength}
\statementsource{Analysis-of-Boolean-Functions-by-Ryan-ODonnell}{
  pdf-ffbe9c7b3306-p183-b003,
  pdf-ffbe9c7b3306-p198-b004
}
\label{FABL.BooleanCircuit.circuitQuadraticPCPPProofLength}
\leanok
\uses{FABL.BooleanCircuit, FABL.BooleanCircuit.circuitQuadraticVariableCount}
Let $C$ be a Boolean circuit with $n$ input bits and $g$ gates, and set $m=1+n+g$. The
associated proof length is $m+2^m+2^{m^2}$.
\end{definition}

\begin{definition}[Extension proof index]
\lean{FABL.BooleanCircuit.circuitQuadraticExtensionProofIndex}
\label{FABL.BooleanCircuit.circuitQuadraticExtensionProofIndex}
\leanok
\uses{FABL.BooleanCircuit, FABL.BooleanCircuit.circuitQuadraticPCPPProofLength, FABL.BooleanCircuit.circuitQuadraticVariableCount}
For a Boolean circuit $C$, let $V$ be its quadratic variable count and let
$P=V+2^V+2^{V^2}$ be its quadratic proof length. Each index $i$ among the $V$
extension-proof positions is sent to position $i$ in the full proof, occupying its
initial block of $V$ positions.
\end{definition}

\begin{definition}[Inner proof index]
\lean{FABL.BooleanCircuit.circuitQuadraticInnerProofIndex}
\label{FABL.BooleanCircuit.circuitQuadraticInnerProofIndex}
\leanok
\uses{FABL.BooleanCircuit, FABL.BooleanCircuit.circuitQuadraticPCPPProofLength, FABL.BooleanCircuit.circuitQuadraticVariableCount}
For a Boolean circuit $C$, let $V$ be its quadratic variable count and let
$P=V+2^V+2^{V^2}$ be its quadratic proof length. Each index $i$ among the $2^V+2^{V^2}$
inner-proof positions is sent to position $V+i$ in the full proof.
\end{definition}

\begin{definition}[Quadratic extension of a circuit proof]
\lean{FABL.BooleanCircuit.circuitQuadraticExtensionProof}
\statementsource{Analysis-of-Boolean-Functions-by-Ryan-ODonnell}{pdf-ffbe9c7b3306-p198-b004}
\label{FABL.BooleanCircuit.circuitQuadraticExtensionProof}
\leanok
\uses{FABL.BooleanCircuit, FABL.BooleanCircuit.circuitQuadraticExtensionProofIndex, FABL.BooleanCircuit.circuitQuadraticPCPPProofLength, FABL.BooleanCircuit.circuitQuadraticVariableCount}
Let $C$ be a Boolean circuit with $n$ input bits, and let $p$ be a Boolean string of
length equal to the circuit’s quadratic PCPP proof length. The quadratic extension of
$p$ is the Boolean string on the circuit’s quadratic variable positions whose value at
each position is the value of $p$ at the corresponding extension-proof index.
\end{definition}

\begin{definition}[Quadratic inner proof]
\lean{FABL.BooleanCircuit.circuitQuadraticInnerProof}
\statementsource{Analysis-of-Boolean-Functions-by-Ryan-ODonnell}{pdf-ffbe9c7b3306-p198-b004}
\label{FABL.BooleanCircuit.circuitQuadraticInnerProof}
\leanok
\uses{FABL.BooleanCircuit, FABL.BooleanCircuit.circuitQuadraticInnerProofIndex, FABL.BooleanCircuit.circuitQuadraticPCPPProofLength, FABL.BooleanCircuit.circuitQuadraticVariableCount}
Let $C$ be a Boolean circuit with $n$ input bits, and let $p$ be a Boolean string of
length equal to the circuit’s quadratic PCPP proof length. The quadratic inner proof is
the Boolean string on the inner-proof positions whose value at each position is the
value of $p$ at the corresponding inner-proof index.
\end{definition}

\begin{definition}[Joined quadratic proof]
\lean{FABL.BooleanCircuit.joinCircuitQuadraticProof}
\label{FABL.BooleanCircuit.joinCircuitQuadraticProof}
\leanok
\uses{FABL.BooleanCircuit, FABL.BooleanCircuit.circuitQuadraticPCPPProofLength, FABL.BooleanCircuit.circuitQuadraticVariableCount}
Let $C$ be a Boolean circuit, let $V$ be its quadratic variable count, and let
$P=V+2^V+2^{V^2}$ be its quadratic proof length. Given a Boolean function on the $V$
extension-proof positions and a Boolean function on the $2^V+2^{V^2}$ inner-proof
positions, their joined proof on the $P$ positions agrees with the extension function on
the initial block and with the inner-proof function on the remaining block.
\end{definition}

\begin{theorem}[Extension proof recovered from joined proof]
\lean{FABL.BooleanCircuit.circuitQuadraticExtensionProof_join}
\label{FABL.BooleanCircuit.circuitQuadraticExtensionProof_join}
\leanok
\uses{FABL.BooleanCircuit, FABL.BooleanCircuit.circuitQuadraticExtensionProof, FABL.BooleanCircuit.circuitQuadraticExtensionProofIndex, FABL.BooleanCircuit.circuitQuadraticVariableCount, FABL.BooleanCircuit.joinCircuitQuadraticProof}
\difficulty{3}
Let $C$ be a Boolean circuit on $n$ input bits. For any Boolean assignment to the
circuit’s extension variables and any Boolean inner proof of length $2^m+2^{m^2}$, where
$m$ is the number of quadratic variables associated with $C$, the extension extracted
from the joined proof is the original extension assignment.
\end{theorem}

\begin{theorem}[Inner proof of a joined circuit proof]
\lean{FABL.BooleanCircuit.circuitQuadraticInnerProof_join}
\label{FABL.BooleanCircuit.circuitQuadraticInnerProof_join}
\leanok
\uses{FABL.BooleanCircuit, FABL.BooleanCircuit.circuitQuadraticInnerProof, FABL.BooleanCircuit.circuitQuadraticInnerProofIndex, FABL.BooleanCircuit.circuitQuadraticVariableCount, FABL.BooleanCircuit.joinCircuitQuadraticProof}
\difficulty{3}
Let $C$ be a Boolean circuit with $n$ input bits and $g$ ordered gates, and set
$m=1+n+g$. Given an extension assignment $e\colon\{1,\ldots,m\}\to\{0,1\}$ and an
inner-proof assignment $q\colon\{1,\ldots,2^m+2^{m^2}\}\to\{0,1\}$, form a proof of
length $m+2^m+2^{m^2}$ by placing $e$ in its first $m$ positions and $q$ in the
remaining positions. Restricting this joined proof to those remaining positions recovers
$q$.
\end{theorem}

\begin{definition}[Encoding a circuit’s quadratic PCPP proof]
\lean{FABL.BooleanCircuit.encodeCircuitQuadraticPCPPProof}
\statementsource{Analysis-of-Boolean-Functions-by-Ryan-ODonnell}{pdf-ffbe9c7b3306-p198-b004}
\label{FABL.BooleanCircuit.encodeCircuitQuadraticPCPPProof}
\leanok
\uses{FABL.BooleanCircuit, FABL.BooleanCircuit.circuitQuadraticPCPPProofLength, FABL.BooleanCircuit.circuitQuadraticVariableCount, FABL.BooleanCircuit.circuitQuadraticWitness, FABL.BooleanCircuit.joinCircuitQuadraticProof, FABL.encodeHomogeneousQuadraticProof, FABL.pcppBoolCubeEquiv}
Let $C$ be a Boolean circuit with $n$ input bits, and let $w$ be an assignment of
Boolean values to its input positions. The encoded quadratic PCPP proof is the proof
string obtained by joining the circuit’s quadratic extension of the quadratic witness
associated with $w$ and the homogeneous quadratic proof encoded from that witness.
\end{definition}

\begin{definition}[Inner verifier address map]
\lean{FABL.BooleanCircuit.circuitQuadraticInnerVerifierAddress}
\label{FABL.BooleanCircuit.circuitQuadraticInnerVerifierAddress}
\leanok
\uses{FABL.BooleanCircuit, FABL.BooleanCircuit.circuitQuadraticExtensionProofIndex, FABL.BooleanCircuit.circuitQuadraticInnerProofIndex, FABL.BooleanCircuit.circuitQuadraticPCPPProofLength, FABL.BooleanCircuit.circuitQuadraticVariableCount}
For a Boolean circuit $C$, let $V$ be its quadratic variable count and let
$M=2^V+2^{V^2}$ be the number of inner-proof positions. The map from the disjoint union
of the $V$ extension positions and the $M$ inner-proof positions to the disjoint union
of the $n$ circuit inputs and the $V+M$ proof positions sends each extension position
$i$ to proof position $i$ and each inner-proof position $j$ to proof position $V+j$.
\end{definition}

\begin{definition}[Quadratic inner-function tester]
\lean{FABL.BooleanCircuit.circuitQuadraticInnerFunctionTester}
\statementsource{Analysis-of-Boolean-Functions-by-Ryan-ODonnell}{pdf-ffbe9c7b3306-p198-b003}
\label{FABL.BooleanCircuit.circuitQuadraticInnerFunctionTester}
\leanok
\uses{FABL.BooleanCircuit, FABL.BooleanCircuit.circuitQuadraticInnerVerifierAddress, FABL.BooleanCircuit.circuitQuadraticPCPPProofLength, FABL.BooleanCircuit.circuitQuadraticSystem, FABL.NonadaptiveFunctionTester, FABL.homogeneousQuadraticPCPPSystem, FABL.reindexDomainFunctionTester}
Let $C$ be a Boolean circuit with $n$ input bits. Its quadratic inner-function tester is
a nonadaptive Boolean-function tester making four queries, with addresses drawn from the
disjoint union of the input positions and the positions of a quadratic PCPP proof; it is
obtained by reindexing the verifier for the homogeneous quadratic PCPP associated with
the circuit’s quadratic system.
\end{definition}

\begin{theorem}[Acceptance probability identity]
\lean{FABL.BooleanCircuit.circuitQuadraticInnerFunctionTester_acceptanceProbability}
\label{FABL.BooleanCircuit.circuitQuadraticInnerFunctionTester_acceptanceProbability}
\leanok
\uses{FABL.BooleanCircuit, FABL.BooleanCircuit.circuitQuadraticExtensionProof, FABL.BooleanCircuit.circuitQuadraticInnerFunctionTester, FABL.BooleanCircuit.circuitQuadraticInnerProof, FABL.BooleanCircuit.circuitQuadraticPCPPProofLength, FABL.BooleanCircuit.circuitQuadraticSystem, FABL.NonadaptiveFunctionTester.acceptanceProbability, FABL.PCPPSystem.acceptanceProbability, FABL.homogeneousQuadraticPCPPSystem, FABL.inputProofWord, FABL.reindexDomainFunctionTester_acceptanceProbability}
\difficulty{4}
Let $C$ be a Boolean circuit on $n$ input bits, let $w$ be an input word, and let $p$ be
a Boolean proof word of the circuit’s prescribed proof length. The inner function
tester’s acceptance probability on the combined word $(w,p)$ equals the acceptance
probability of the homogeneous quadratic PCPP verifier for the circuit’s quadratic
system, using the extension proof and inner proof obtained by restricting $p$ to their
respective portions.
\end{theorem}

\begin{definition}[Quadratic linear proof index]
\lean{FABL.BooleanCircuit.circuitQuadraticLinearProofIndex}
\label{FABL.BooleanCircuit.circuitQuadraticLinearProofIndex}
\leanok
\uses{FABL.BooleanCircuit, FABL.BooleanCircuit.circuitQuadraticInnerProofIndex, FABL.BooleanCircuit.circuitQuadraticPCPPProofLength, FABL.BooleanCircuit.circuitQuadraticVariableCount, FABL.quadraticLinearProofIndex}
For a Boolean circuit $C$ with $n$ input bits and $g$ gates, let $q=1+n+g$ and let the
proof length be $q+2^q+2^{q^2}$. The quadratic linear proof index assigns each index in
$\text{Fin}(2^q)$ an index in $\text{Fin}(q+2^q+2^{q^2})$ by first applying the quadratic
linear proof-index map and then the circuit's inner proof-index map.
\end{definition}

\begin{definition}[Three-query input-consistency tester]
\lean{FABL.BooleanCircuit.circuitQuadraticInputConsistencyThreeTester}
\statementsource{Analysis-of-Boolean-Functions-by-Ryan-ODonnell}{pdf-ffbe9c7b3306-p198-b004}
\label{FABL.BooleanCircuit.circuitQuadraticInputConsistencyThreeTester}
\leanok
\uses{FABL.BooleanCircuit, FABL.BooleanCircuit.circuitQuadraticInputIndex, FABL.BooleanCircuit.circuitQuadraticLinearProofIndex, FABL.BooleanCircuit.circuitQuadraticPCPPProofLength, FABL.BooleanCircuit.circuitQuadraticVariableCount, FABL.F₂Cube, FABL.NonadaptiveFunctionTester, FABL.affineConsistencyTranscriptAccept, FABL.f₂CubeIndexEquiv, FABL.f₂StandardBasis, FABL.uniformPMF}
Let $C$ be a Boolean circuit with $n>0$ input bits. Its input-consistency tester makes
three nonadaptive Boolean queries, with addresses in the disjoint union of input
positions and quadratic PCPP proof positions. It chooses uniformly a pair consisting of
an input position $i$ and a point $y$ in the Boolean cube of dimension equal to the
circuit’s quadratic variable count; it queries input position $i$ and the two
linear-proof coordinates corresponding to $y$ and to $y$ shifted by the standard basis
vector for input position $i$, and accepts according to the affine-consistency
predicate.
\end{definition}

\begin{definition}[Four-query input-consistency tester]
\lean{FABL.BooleanCircuit.circuitQuadraticInputConsistencyFunctionTester}
\statementsource{Analysis-of-Boolean-Functions-by-Ryan-ODonnell}{pdf-ffbe9c7b3306-p198-b004}
\label{FABL.BooleanCircuit.circuitQuadraticInputConsistencyFunctionTester}
\leanok
\uses{FABL.BooleanCircuit, FABL.BooleanCircuit.circuitQuadraticInputConsistencyThreeTester, FABL.BooleanCircuit.circuitQuadraticPCPPProofLength, FABL.NonadaptiveFunctionTester, FABL.padThreeQueryFunctionTester}
Let $C$ be a Boolean circuit with $n>0$ input bits. Its four-query input-consistency
tester is obtained by padding the three-query input-consistency tester with one query,
and has addresses in the disjoint union of input positions and quadratic PCPP proof
positions.
\end{definition}

\begin{theorem}[Input consistency acceptance probability]
\lean{FABL.BooleanCircuit.circuitQuadraticInputConsistencyFunctionTester_acceptanceProbability}
\label{FABL.BooleanCircuit.circuitQuadraticInputConsistencyFunctionTester_acceptanceProbability}
\leanok
\uses{FABL.BooleanCircuit, FABL.BooleanCircuit.circuitQuadraticInnerProof, FABL.BooleanCircuit.circuitQuadraticInputConsistencyFunctionTester, FABL.BooleanCircuit.circuitQuadraticInputConsistencyThreeTester, FABL.BooleanCircuit.circuitQuadraticInputIndex, FABL.BooleanCircuit.circuitQuadraticPCPPProofLength, FABL.NonadaptiveFunctionTester.acceptanceProbability, FABL.NonadaptiveFunctionTester.transcript_apply, FABL.affineConsistencyTranscriptAccept, FABL.affineCoordinateCorrectionProbability, FABL.f₂CubeProofFunction, FABL.inputProofWord, FABL.padThreeQueryFunctionTester_acceptanceProbability, FABL.pmfExpectation_uniformPMF_eq_expect, FABL.quadraticLinearProofTable, FABL.uniformPMF, FABL.uniformProbability}
\difficulty{4}
Let $n>0$, let $C$ be a Boolean circuit on $n$ input bits, let $w$ be an input word, and
let $p$ be a Boolean proof word of the circuit’s prescribed length. The input
consistency tester’s acceptance probability on $(w,p)$ equals the average over input
coordinates $i$ of the probability, for uniformly chosen $y\in\mathbb{F}_2^n$, that the
local correction of the sign-valued function encoded by the linear portion of the inner
proof agrees with the sign encoding of $w_i$.
\end{theorem}

\begin{theorem}[Mismatch and affine-correction rejection]
\lean{FABL.BooleanCircuit.half_bool_mismatch_le_mappedAffineCorrection_rejection}
\label{FABL.BooleanCircuit.half_bool_mismatch_le_mappedAffineCorrection_rejection}
\leanok
\uses{FABL.F₂Cube, FABL.Sign, FABL.affineCoordinateCorrectionProbability, FABL.affineCoordinateCorrectionProbability_le_one, FABL.f₂StandardBasis, FABL.localCorrection, FABL.uniformProbability, FABL.uniformProbability_sign_eq_one_add_neg_one_eq_one}
\difficulty{5}
Let $f:\mathbb F_2^N\to\{-1,1\}$, let $i$ be a coordinate, and let $b,b_0$ be Boolean
values. If the affine-coordinate correction associated with $f$ and $i$ returns the sign
encoding of $b_0$ with probability at least $1/2$ under a uniform choice of its
auxiliary point, then $\frac12$ when $b\ne b_0$ (and $0$ when $b=b_0$) is at most one
minus the probability that the same correction returns the sign encoding of $b$.
\end{theorem}

\begin{theorem}[Input distance and consistency-test rejection]
\lean{FABL.BooleanCircuit.half_relativeHammingDist_le_circuitQuadraticInputConsistency_rejection}
\label{FABL.BooleanCircuit.half_relativeHammingDist_le_circuitQuadraticInputConsistency_rejection}
\leanok
\uses{FABL.BooleanCircuit, FABL.BooleanCircuit.circuitQuadraticInnerProof, FABL.BooleanCircuit.circuitQuadraticInputConsistencyFunctionTester, FABL.BooleanCircuit.circuitQuadraticInputConsistencyFunctionTester_acceptanceProbability, FABL.BooleanCircuit.circuitQuadraticInputIndex, FABL.BooleanCircuit.circuitQuadraticPCPPProofLength, FABL.BooleanCircuit.half_bool_mismatch_le_mappedAffineCorrection_rejection, FABL.NonadaptiveFunctionTester.acceptanceProbability, FABL.affineCoordinateCorrectionProbability, FABL.f₂CubeProofFunction, FABL.inputProofWord, FABL.quadraticLinearProofTable, FABL.relativeHammingDist, FABL.uniformProbability, FABL.uniformProbability_ne_eq_relativeHammingDist}
\difficulty{4}
Let $C$ be a Boolean circuit on $n$ input bits, where $n>0$, let $p$ be a Boolean proof
word of the prescribed length for $C$, and let $w,w_0\in\{0,1\}^n$. Suppose that for
every input coordinate, the affine-coordinate correction derived from the quadratic
proof table of $p$ returns the sign encoding of the corresponding bit of $w_0$ with
probability at least $1/2$. Then one half the relative Hamming distance between $w$ and
$w_0$ is at most one minus the acceptance probability of the circuit's input-consistency
tester on the concatenation of $w$ and $p$.
\end{theorem}

\begin{theorem}[Circuit quadratic inner function tester complete]
\lean{FABL.BooleanCircuit.circuitQuadraticInnerFunctionTester_complete}
\statementsource{Analysis-of-Boolean-Functions-by-Ryan-ODonnell}{pdf-ffbe9c7b3306-p198-b003}
\label{FABL.BooleanCircuit.circuitQuadraticInnerFunctionTester_complete}
\leanok
\uses{FABL.BooleanCircuit, FABL.BooleanCircuit.acceptedProperty, FABL.BooleanCircuit.circuitQuadraticExtensionProof_join, FABL.BooleanCircuit.circuitQuadraticInnerFunctionTester, FABL.BooleanCircuit.circuitQuadraticInnerFunctionTester_acceptanceProbability, FABL.BooleanCircuit.circuitQuadraticInnerProof_join, FABL.BooleanCircuit.circuitQuadraticSystem, FABL.BooleanCircuit.circuitQuadraticVariableCount, FABL.BooleanCircuit.circuitQuadraticWitness, FABL.BooleanCircuit.circuitQuadraticWitness_satisfies, FABL.BooleanCircuit.encodeCircuitQuadraticPCPPProof, FABL.IsBoolHomogeneousQuadraticSystem, FABL.NonadaptiveFunctionTester.acceptanceProbability, FABL.SatisfiesHomogeneousQuadraticSystem, FABL.homogeneousQuadraticPCPPSystem_complete, FABL.inputProofWord, FABL.pcppBoolCubeEquiv}
\difficulty{4}
% dataset-pipeline: stub -- prose pending build_informal_statements.py
Statement pending regeneration.
\end{theorem}

\begin{theorem}[Completeness of the input consistency tester]
\lean{FABL.BooleanCircuit.circuitQuadraticInputConsistencyFunctionTester_complete}
\proofsource{Analysis-of-Boolean-Functions-by-Ryan-ODonnell}{pdf-ffbe9c7b3306-p198-b004}
\label{FABL.BooleanCircuit.circuitQuadraticInputConsistencyFunctionTester_complete}
\leanok
\uses{FABL.BooleanCircuit, FABL.BooleanCircuit.circuitQuadraticInnerProof_join, FABL.BooleanCircuit.circuitQuadraticInputConsistencyFunctionTester, FABL.BooleanCircuit.circuitQuadraticInputConsistencyFunctionTester_acceptanceProbability, FABL.BooleanCircuit.circuitQuadraticInputIndex, FABL.BooleanCircuit.circuitQuadraticVariableCount, FABL.BooleanCircuit.circuitQuadraticWitness, FABL.BooleanCircuit.circuitQuadraticWitness_input, FABL.BooleanCircuit.encodeCircuitQuadraticPCPPProof, FABL.F₂Cube, FABL.NonadaptiveFunctionTester.acceptanceProbability, FABL.affineCoordinateCorrectionProbability, FABL.f₂CubeProofFunction_quadraticLinearProofTable, FABL.f₂DotProduct, FABL.f₂SignCharacter, FABL.f₂SignCharacter_standardBasis, FABL.f₂StandardBasis, FABL.inputProofWord, FABL.localCorrection, FABL.localCorrection_f₂SignCharacter, FABL.pcppBoolSignEquiv, FABL.pcppBoolSignEquiv_eq_signEncode, FABL.quadraticLinearProofFunction_encodeHomogeneousQuadraticProof, FABL.signEncode, FABL.uniformProbability}
\difficulty{4}
Let $n>0$, let $C$ be a Boolean circuit on $n$ input bits, and let $w$ be any input
word. When given $w$ together with the proof encoding the circuit’s canonical quadratic
witness for $w$, the circuit’s input consistency tester accepts with probability $1$.
\end{theorem}

\begin{definition}[Four-query circuit tester]
\lean{FABL.BooleanCircuit.circuitQuadraticFourQueryFunctionTester}
\statementsource{Analysis-of-Boolean-Functions-by-Ryan-ODonnell}{pdf-ffbe9c7b3306-p198-b003}
\label{FABL.BooleanCircuit.circuitQuadraticFourQueryFunctionTester}
\leanok
\uses{FABL.BooleanCircuit, FABL.BooleanCircuit.circuitQuadraticInnerFunctionTester, FABL.BooleanCircuit.circuitQuadraticInputConsistencyFunctionTester, FABL.BooleanCircuit.circuitQuadraticPCPPProofLength, FABL.NonadaptiveFunctionTester, FABL.uniformMixtureFunctionTester}
Let $C$ be a Boolean circuit with $n>0$ input bits. Its four-query function tester is
the uniform mixture of its quadratic inner-function tester and its four-query
input-consistency tester; it has Boolean answers and addresses in the disjoint union of
input positions and quadratic PCPP proof positions.
\end{definition}

\begin{definition}[Four-query circuit PCPP system]
\lean{FABL.BooleanCircuit.circuitQuadraticFourQueryPCPPSystem}
\statementsource{Analysis-of-Boolean-Functions-by-Ryan-ODonnell}{pdf-ffbe9c7b3306-p198-b003}
\label{FABL.BooleanCircuit.circuitQuadraticFourQueryPCPPSystem}
\leanok
\uses{FABL.BooleanCircuit, FABL.BooleanCircuit.circuitQuadraticFourQueryFunctionTester, FABL.BooleanCircuit.circuitQuadraticPCPPProofLength, FABL.PCPPSystem}
Let $C$ be a Boolean circuit with $n>0$ input bits. The four-query circuit PCPP system
has Boolean alphabet, input length $n$, proof length equal to the circuit’s quadratic
PCPP proof length, and verifier given by the circuit’s four-query function tester.
\end{definition}

\begin{theorem}[Acceptance probability of the combined verifier]
\lean{FABL.BooleanCircuit.circuitQuadraticFourQueryPCPPSystem_acceptanceProbability}
\label{FABL.BooleanCircuit.circuitQuadraticFourQueryPCPPSystem_acceptanceProbability}
\leanok
\uses{FABL.BooleanCircuit, FABL.BooleanCircuit.circuitQuadraticFourQueryFunctionTester, FABL.BooleanCircuit.circuitQuadraticFourQueryPCPPSystem, FABL.BooleanCircuit.circuitQuadraticInnerFunctionTester, FABL.BooleanCircuit.circuitQuadraticInputConsistencyFunctionTester, FABL.BooleanCircuit.circuitQuadraticPCPPProofLength, FABL.NonadaptiveFunctionTester.acceptanceProbability, FABL.PCPPSystem.acceptanceProbability, FABL.inputProofWord, FABL.uniformMixtureFunctionTester_acceptanceProbability}
\difficulty{3}
Let $C$ be a Boolean circuit on $n$ input bits, with $n>0$. For every Boolean input $w$
and every Boolean proof of the circuit’s quadratic PCPP proof length, the combined
verifier’s acceptance probability equals the average of the acceptance probabilities of
its inner-function tester and input-consistency tester, each applied to the word formed
from $w$ and the proof.
\end{theorem}

\begin{theorem}[Completeness of the combined verifier]
\lean{FABL.BooleanCircuit.circuitQuadraticFourQueryPCPPSystem_complete}
\label{FABL.BooleanCircuit.circuitQuadraticFourQueryPCPPSystem_complete}
\leanok
\uses{FABL.BooleanCircuit, FABL.BooleanCircuit.acceptedProperty, FABL.BooleanCircuit.circuitQuadraticFourQueryPCPPSystem, FABL.BooleanCircuit.circuitQuadraticFourQueryPCPPSystem_acceptanceProbability, FABL.BooleanCircuit.circuitQuadraticInnerFunctionTester_complete, FABL.BooleanCircuit.circuitQuadraticInputConsistencyFunctionTester_complete, FABL.BooleanCircuit.encodeCircuitQuadraticPCPPProof, FABL.PCPPSystem.acceptanceProbability}
\difficulty{3}
Let $C$ be a Boolean circuit on $n$ input bits, with $n>0$. If a Boolean input $w$ is
accepted by $C$, then the combined four-query verifier accepts $w$ with probability $1$
when given the encoded circuit quadratic proof for $w$.
\end{theorem}

\begin{definition}[Quadratic PCPP rejection rate]
\lean{FABL.BooleanCircuit.circuitQuadraticPCPPRejectionRate}
\statementsource{Analysis-of-Boolean-Functions-by-Ryan-ODonnell}{pdf-ffbe9c7b3306-p198-b003}
\label{FABL.BooleanCircuit.circuitQuadraticPCPPRejectionRate}
\leanok
The rejection rate assigned to the quadratic circuit PCPP construction is the constant
$1/200000$.
\end{definition}

\begin{theorem}[Component acceptance from combined acceptance]
\lean{FABL.BooleanCircuit.circuitQuadraticFourQueryPCPPSystem_component_acceptanceProbability_ge}
\label{FABL.BooleanCircuit.circuitQuadraticFourQueryPCPPSystem_component_acceptanceProbability_ge}
\leanok
\uses{FABL.BooleanCircuit, FABL.BooleanCircuit.circuitQuadraticFourQueryFunctionTester, FABL.BooleanCircuit.circuitQuadraticFourQueryPCPPSystem, FABL.BooleanCircuit.circuitQuadraticInnerFunctionTester, FABL.BooleanCircuit.circuitQuadraticInputConsistencyFunctionTester, FABL.BooleanCircuit.circuitQuadraticPCPPProofLength, FABL.BooleanCircuit.circuitQuadraticPCPPRejectionRate, FABL.NonadaptiveFunctionTester, FABL.NonadaptiveFunctionTester.acceptanceProbability, FABL.PCPPSystem.acceptanceProbability, FABL.component_acceptanceProbability_ge_of_uniformMixture_ge, FABL.homogeneousQuadraticPCPPRejectionRate, FABL.inputProofWord, FABL.uniformMixtureFunctionTester}
\difficulty{4}
Let $C$ be a Boolean circuit on $n$ input bits, with $n>0$, let $psilon$ be real, and
let $w$ and $p$ be a Boolean input and a Boolean proof of the circuit’s quadratic PCPP
proof length. If the combined verifier accepts with probability at least
$1-
ho_Cpsilon$, where $
ho_C$ is the circuit quadratic rejection rate, then each of
its inner-function and input-consistency testers accepts the combined input-proof word
with probability at least $1-
ho_Qepsilon$, where $
ho_Q$ is the homogeneous quadratic
PCPP rejection rate.
\end{theorem}

\begin{theorem}[Circuit quadratic four query pcpp system close of acceptance probability ge]
\lean{FABL.BooleanCircuit.circuitQuadraticFourQueryPCPPSystem_close_of_acceptanceProbability_ge}
\statementsource{Analysis-of-Boolean-Functions-by-Ryan-ODonnell}{pdf-ffbe9c7b3306-p198-b003}
\label{FABL.BooleanCircuit.circuitQuadraticFourQueryPCPPSystem_close_of_acceptanceProbability_ge}
\leanok
\uses{FABL.BooleanCircuit, FABL.BooleanCircuit.acceptedProperty, FABL.BooleanCircuit.circuitQuadraticExtensionProof, FABL.BooleanCircuit.circuitQuadraticFourQueryPCPPSystem, FABL.BooleanCircuit.circuitQuadraticFourQueryPCPPSystem_component_acceptanceProbability_ge, FABL.BooleanCircuit.circuitQuadraticInnerFunctionTester_acceptanceProbability, FABL.BooleanCircuit.circuitQuadraticInnerProof, FABL.BooleanCircuit.circuitQuadraticInput, FABL.BooleanCircuit.circuitQuadraticInputConsistencyFunctionTester, FABL.BooleanCircuit.circuitQuadraticInputIndex, FABL.BooleanCircuit.circuitQuadraticInput_accepted_of_satisfies, FABL.BooleanCircuit.circuitQuadraticPCPPProofLength, FABL.BooleanCircuit.circuitQuadraticPCPPRejectionRate, FABL.BooleanCircuit.circuitQuadraticSystem, FABL.BooleanCircuit.circuitQuadraticVariableCount, FABL.BooleanCircuit.half_relativeHammingDist_le_circuitQuadraticInputConsistency_rejection, FABL.IsClose, FABL.IsCloseToProperty, FABL.NonadaptiveFunctionTester.acceptanceProbability, FABL.PCPPSystem.acceptanceProbability, FABL.affineCoordinateCorrectionProbability, FABL.f₂CubeProofFunction, FABL.homogeneousQuadraticPCPPRejectionRate, FABL.homogeneousQuadraticPCPPSystem_character_witness_of_acceptanceProbability_ge, FABL.inputProofWord, FABL.pcppBoolCubeEquiv, FABL.quadraticLinearProofTable, FABL.quadraticLinearProof_localCorrectionProbability_ge}
\difficulty{5}
% dataset-pipeline: stub -- prose pending build_informal_statements.py
Statement pending regeneration.
\end{theorem}

\begin{theorem}[Circuit quadratic four query pcpp system is pcpp]
\lean{FABL.BooleanCircuit.circuitQuadraticFourQueryPCPPSystem_isPCPP}
\statementsource{Analysis-of-Boolean-Functions-by-Ryan-ODonnell}{pdf-ffbe9c7b3306-p198-b003}
\label{FABL.BooleanCircuit.circuitQuadraticFourQueryPCPPSystem_isPCPP}
\leanok
\uses{FABL.BooleanCircuit, FABL.BooleanCircuit.acceptedProperty, FABL.BooleanCircuit.circuitQuadraticFourQueryPCPPSystem, FABL.BooleanCircuit.circuitQuadraticFourQueryPCPPSystem_close_of_acceptanceProbability_ge, FABL.BooleanCircuit.circuitQuadraticFourQueryPCPPSystem_complete, FABL.BooleanCircuit.circuitQuadraticPCPPRejectionRate, FABL.BooleanCircuit.encodeCircuitQuadraticPCPPProof, FABL.IsPCPP, FABL.isPCPP_of_exists_close_of_acceptanceProbability_ge}
\difficulty{3}
% dataset-pipeline: stub -- prose pending build_informal_statements.py
Statement pending regeneration.
\end{theorem}

\begin{definition}[Four-query circuit PCPP description]
\lean{FABL.BooleanCircuit.circuitQuadraticFourQueryPCPPDescription}
\label{FABL.BooleanCircuit.circuitQuadraticFourQueryPCPPDescription}
\leanok
\uses{FABL.BooleanCircuit, FABL.BooleanCircuit.acceptedProperty, FABL.BooleanCircuit.circuitQuadraticFourQueryPCPPSystem, FABL.BooleanCircuit.circuitQuadraticFourQueryPCPPSystem_isPCPP, FABL.BooleanCircuit.circuitQuadraticPCPPProofLength, FABL.BooleanCircuit.circuitQuadraticPCPPRejectionRate, FABL.PCPPDescription}
For a Boolean circuit $C$ with $n$ input bits, where $n$ is nonzero, the four-query
circuit PCPP description specifies a PCPP for the property that $C$ accepts its input.
Its proof length is $q+2^q+2^{q^2}$, where $q=1+n+g$ and $g$ is the number of gates of
$C$; its verifier makes four queries and has the circuit quadratic PCPP rejection rate.
\end{definition}

\begin{theorem}[Circuit quadratic four query pcpp description proof length]
\lean{FABL.BooleanCircuit.circuitQuadraticFourQueryPCPPDescription_proofLength}
\statementsource{Analysis-of-Boolean-Functions-by-Ryan-ODonnell}{pdf-ffbe9c7b3306-p198-b003}
\label{FABL.BooleanCircuit.circuitQuadraticFourQueryPCPPDescription_proofLength}
\leanok
\uses{FABL.BooleanCircuit, FABL.BooleanCircuit.circuitQuadraticFourQueryPCPPDescription, FABL.BooleanCircuit.circuitQuadraticVariableCount}
\difficulty{1}
% dataset-pipeline: stub -- prose pending build_informal_statements.py
Statement pending regeneration.
\end{theorem}

\begin{theorem}[Quadratic PCPP proof length bound]
\lean{FABL.BooleanCircuit.circuitQuadraticPCPPProofLength_le_three_mul_two_pow_square}
\label{FABL.BooleanCircuit.circuitQuadraticPCPPProofLength_le_three_mul_two_pow_square}
\leanok
\uses{FABL.BooleanCircuit, FABL.BooleanCircuit.circuitQuadraticPCPPProofLength, FABL.BooleanCircuit.circuitQuadraticVariableCount}
\difficulty{4}
For every Boolean circuit $C$ with $n$ input bits and $g$ gates, its quadratic PCPP
proof length is at most $3\cdot 2^{(n+g+1)^2}$.
\end{theorem}

\begin{theorem}[Acceptance of a circuit with no input bits]
\lean{FABL.BooleanCircuit.acceptedProperty_zero_universal_or_empty}
\label{FABL.BooleanCircuit.acceptedProperty_zero_universal_or_empty}
\leanok
\uses{FABL.BooleanCircuit, FABL.BooleanCircuit.acceptedProperty}
\difficulty{4}
Let $C$ be a Boolean circuit with zero input bits. Either every input assignment is
accepted by $C$, or no input assignment is accepted by $C$.
\end{theorem}

\begin{theorem}[Four-query PCPP description or degenerate case]
\lean{FABL.BooleanCircuit.circuitQuadraticFourQueryPCPPDescription_or_zero}
\label{FABL.BooleanCircuit.circuitQuadraticFourQueryPCPPDescription_or_zero}
\leanok
\uses{FABL.BooleanCircuit, FABL.BooleanCircuit.acceptedProperty, FABL.BooleanCircuit.acceptedProperty_zero_universal_or_empty, FABL.BooleanCircuit.circuitQuadraticFourQueryPCPPDescription, FABL.PCPPDescription}
\difficulty{4}
Let $C$ be a Boolean circuit on $n$ input bits. Either $n=0$ and the circuit accepts
every input or rejects every input, or, for some positive $n$, the property of inputs
accepted by $C$ has a nonadaptive probabilistically checkable proof of proximity
description.
\end{theorem}

\begin{definition}[Three-query circuit PCPP description]
\lean{FABL.BooleanCircuit.circuitQuadraticThreeQueryPCPPDescription}
\statementsource{Analysis-of-Boolean-Functions-by-Ryan-ODonnell}{
  pdf-ffbe9c7b3306-p183-b003,
  pdf-ffbe9c7b3306-p198-b004
}
\label{FABL.BooleanCircuit.circuitQuadraticThreeQueryPCPPDescription}
\leanok
\uses{FABL.BooleanCircuit, FABL.BooleanCircuit.acceptedProperty, FABL.BooleanCircuit.circuitQuadraticFourQueryPCPPDescription, FABL.PCPPDescription, FABL.threeQueryPCPPDescription}
For a Boolean circuit $C$ on $n$ input bits, with $n>0$, the quadratic circuit
construction assigns a PCPP description for the property that the circuit accepts its
input. The description has three queries and satisfies the perfect-completeness and
rejection-rate conditions of a PCPP for that property.
\end{definition}

\begin{theorem}[Three-query PCPP query count]
\lean{FABL.BooleanCircuit.circuitQuadraticThreeQueryPCPPDescription_queryCount}
\proofsource{Analysis-of-Boolean-Functions-by-Ryan-ODonnell}{pdf-ffbe9c7b3306-p183-b003}
\label{FABL.BooleanCircuit.circuitQuadraticThreeQueryPCPPDescription_queryCount}
\leanok
\uses{FABL.BooleanCircuit, FABL.BooleanCircuit.circuitQuadraticThreeQueryPCPPDescription}
\difficulty{1}
Let $C$ be a Boolean circuit with $n>0$ input bits. Its quadratic three-query PCPP
description has query count exactly $3$.
\end{theorem}

\begin{theorem}[Circuit quadratic three query pcpp description rejection rate]
\lean{FABL.BooleanCircuit.circuitQuadraticThreeQueryPCPPDescription_rejectionRate}
\statementsource{Analysis-of-Boolean-Functions-by-Ryan-ODonnell}{pdf-ffbe9c7b3306-p183-b003}
\label{FABL.BooleanCircuit.circuitQuadraticThreeQueryPCPPDescription_rejectionRate}
\leanok
\uses{FABL.BooleanCircuit, FABL.BooleanCircuit.circuitQuadraticFourQueryPCPPDescription, FABL.BooleanCircuit.circuitQuadraticPCPPRejectionRate, FABL.BooleanCircuit.circuitQuadraticThreeQueryPCPPDescription, FABL.threeQueryPCPPDescription_rejectionRate}
\difficulty{3}
% dataset-pipeline: stub -- prose pending build_informal_statements.py
Statement pending regeneration.
\end{theorem}

\begin{theorem}[Three-query PCPP proof length]
\lean{FABL.BooleanCircuit.circuitQuadraticThreeQueryPCPPDescription_proofLength}
\label{FABL.BooleanCircuit.circuitQuadraticThreeQueryPCPPDescription_proofLength}
\leanok
\uses{FABL.BooleanCircuit, FABL.BooleanCircuit.circuitQuadraticFourQueryPCPPSystem, FABL.BooleanCircuit.circuitQuadraticPCPPProofLength, FABL.BooleanCircuit.circuitQuadraticThreeQueryPCPPDescription, FABL.PCPPSystem.descriptionSize, FABL.threeQueryPCPPDescription_proofLength}
\difficulty{3}
Let $C$ be a Boolean circuit with $n>0$ input bits. The proof length of its three-query
PCPP description is the sum of the quadratic PCPP proof length for $C$ and $64$ times
the description size of the associated four-query PCPP system.
\end{theorem}

\begin{theorem}[Split Character Seed Cardinality Bound]
\lean{FABL.splitCharacterSeed_card_le_three_cube_powers}
\label{FABL.splitCharacterSeed_card_le_three_cube_powers}
\leanok
\uses{FABL.BinarySubspaceSplit, FABL.BinarySubspaceSplit.point, FABL.F₂Cube, FABL.SplitCharacterSeed, FABL.𝔽₂}
\difficulty{5}
Let $d$ be a natural number, let $H$ be a subspace of $\bbf_2^d$, and let a binary
subspace split of $H$ have complementary dimension $k$. The number of seeds
$(x,y,h)\in\bbf_2^k\times\bbf_2^k\times H$ is at most $2^{3d}$.
\end{theorem}

\begin{theorem}[Seed-space cardinality bound]
\lean{FABL.linearizedQuadraticSystemCharacterTester_seed_card_le}
\label{FABL.linearizedQuadraticSystemCharacterTester_seed_card_le}
\leanok
\uses{FABL.HomogeneousQuadraticSystem, FABL.SatisfiesLinearizedQuadraticSystem, FABL.SplitCharacterSeed, FABL.binarySubspaceSplit, FABL.homogeneousQuadraticSystemDirection, FABL.linearizedQuadraticSystemCharacterTester, FABL.perpendicularSubspace, FABL.splitCharacterSeed_card_le_three_cube_powers}
\difficulty{4}
For every natural number $d$ and every finite system $L$ of homogeneous quadratic
equations in dimension $d$, the seed space of the linearized-system character tester for
$L$ has cardinality at most $2^{3d^2}$.
\end{theorem}

\begin{theorem}[Seed count of the quadratic character pair tester]
\lean{FABL.quadraticCharacterPairTester_seed_card_le}
\label{FABL.quadraticCharacterPairTester_seed_card_le}
\leanok
\uses{FABL.F₂Cube, FABL.quadraticCharacterPairTester, FABL.quadraticPairLinearBLRTester, FABL.quadraticPairMatrixBLRTester, FABL.quadraticPairRankOneTester}
\difficulty{5}
For every natural number $d$, the finite seed space of the quadratic character pair
tester on dimension $d$ has cardinality at most $3\cdot 2^{3(d+1)^2}$.
\end{theorem}

\begin{theorem}[Affine tester seed-space bound]
\lean{FABL.affineConsistencyFunctionTester_seed_card_le}
\label{FABL.affineConsistencyFunctionTester_seed_card_le}
\leanok
\uses{FABL.F₂Cube, FABL.affineConsistencyFunctionTester}
\difficulty{4}
For every positive natural number $d$, the seed space of the affine consistency tester
on dimension $d$ has cardinality at most $2^{3(d+1)^2}$.
\end{theorem}

\begin{theorem}[Seed Bound for the Quadratic Pair Tester]
\lean{FABL.quadraticPairProofFunctionTester_seed_card_le}
\label{FABL.quadraticPairProofFunctionTester_seed_card_le}
\leanok
\uses{FABL.quadraticCharacterPairTester, FABL.quadraticCharacterPairTester_seed_card_le, FABL.quadraticPairProofFunctionTester}
\difficulty{1}
For every natural number $d$, the seed space of the Boolean-valued tester obtained by
combining the linearity, matrix-linearity, and rank-one tests on quadratic-character
addresses has cardinality at most $3\cdot 2^{3(d+1)^2}$.
\end{theorem}

\begin{theorem}[Bound on proof tester seed count]
\lean{FABL.quadraticSystemProofFunctionTester_seed_card_le}
\label{FABL.quadraticSystemProofFunctionTester_seed_card_le}
\leanok
\uses{FABL.HomogeneousQuadraticSystem, FABL.linearizedQuadraticSystemCharacterTester, FABL.linearizedQuadraticSystemCharacterTester_seed_card_le, FABL.quadraticSystemProofFunctionTester}
\difficulty{1}
For every natural number $d$ and every list $L$ of homogeneous quadratic equations in
$d$ variables over $\mathbb F_2$, the seed space of the associated quadratic-system
proof tester has cardinality at most $2^{3(d+1)^2}$.
\end{theorem}

\begin{theorem}[Quadratic consistency tester seed bound]
\lean{FABL.quadraticInputConsistencyFunctionTester_seed_card_le}
\label{FABL.quadraticInputConsistencyFunctionTester_seed_card_le}
\leanok
\uses{FABL.affineConsistencyFunctionTester, FABL.affineConsistencyFunctionTester_seed_card_le, FABL.quadraticInputConsistencyFunctionTester}
\difficulty{1}
For every positive integer $d$, the seed space of the quadratic input consistency tester
with parameter $d$ has cardinality at most $2^{3(d+1)^2}$.
\end{theorem}

\begin{theorem}[Proof system description size bound]
\lean{FABL.homogeneousQuadraticPCPPSystem_descriptionSize_le}
\label{FABL.homogeneousQuadraticPCPPSystem_descriptionSize_le}
\leanok
\uses{FABL.HomogeneousQuadraticSystem, FABL.PCPPSystem.descriptionSize, FABL.homogeneousQuadraticPCPPSystem, FABL.homogeneousQuadraticProofFunctionTester, FABL.quadraticInputConsistencyFunctionTester, FABL.quadraticInputConsistencyFunctionTester_seed_card_le, FABL.quadraticPairProofFunctionTester, FABL.quadraticPairProofFunctionTester_seed_card_le, FABL.quadraticSystemProofFunctionTester, FABL.quadraticSystemProofFunctionTester_seed_card_le}
\difficulty{4}
Let $d>0$ and let $L$ be a finite system of homogeneous quadratic equations in $d$
variables over $\mathbb{F}_2$. The description size of its associated proof system is at
most $5\cdot 2^{3(d+1)^2}$.
\end{theorem}

\begin{theorem}[Bound on the tester seed space]
\lean{FABL.BooleanCircuit.circuitQuadraticInputConsistencyFunctionTester_seed_card_le}
\label{FABL.BooleanCircuit.circuitQuadraticInputConsistencyFunctionTester_seed_card_le}
\leanok
\uses{FABL.BooleanCircuit, FABL.BooleanCircuit.circuitQuadraticInputConsistencyFunctionTester, FABL.BooleanCircuit.circuitQuadraticVariableCount, FABL.F₂Cube}
\difficulty{4}
Let $n>0$ and let $C$ be a Boolean circuit on $n$ input bits. The number of seeds in the
circuit’s input consistency tester is at most $2^{3(q+1)^2}$, where $q=1+n+g$ and $g$ is
the number of gates of $C$.
\end{theorem}

\begin{theorem}[Description size bound for the verifier]
\lean{FABL.BooleanCircuit.circuitQuadraticFourQueryPCPPSystem_descriptionSize_le}
\label{FABL.BooleanCircuit.circuitQuadraticFourQueryPCPPSystem_descriptionSize_le}
\leanok
\uses{FABL.BooleanCircuit, FABL.BooleanCircuit.circuitQuadraticFourQueryFunctionTester, FABL.BooleanCircuit.circuitQuadraticFourQueryPCPPSystem, FABL.BooleanCircuit.circuitQuadraticInnerFunctionTester, FABL.BooleanCircuit.circuitQuadraticInputConsistencyFunctionTester, FABL.BooleanCircuit.circuitQuadraticInputConsistencyFunctionTester_seed_card_le, FABL.BooleanCircuit.circuitQuadraticSystem, FABL.BooleanCircuit.circuitQuadraticVariableCount, FABL.PCPPSystem.descriptionSize, FABL.homogeneousQuadraticPCPPSystem, FABL.homogeneousQuadraticPCPPSystem_descriptionSize_le}
\difficulty{4}
Let $C$ be a Boolean circuit on $n$ input bits, with $n>0$, and let $m$ be its number of
associated quadratic variables. The description size of its four-query PCPP system is at
most $6\cdot 2^{3(m+1)^2}$.
\end{theorem}

\begin{definition}[Four-query circuit PCPP reduction]
\lean{FABL.circuitQuadraticFourQuerySemanticPCPPReduction}
\statementsource{Analysis-of-Boolean-Functions-by-Ryan-ODonnell}{pdf-ffbe9c7b3306-p198-b003}
\label{FABL.circuitQuadraticFourQuerySemanticPCPPReduction}
\leanok
\uses{FABL.BooleanCircuit.circuitQuadraticFourQueryPCPPDescription, FABL.SemanticPCPPReduction}
For every $n>0$ and every Boolean circuit $C$ on $n$ input bits, the reduction assigns a
PCPP description for the property that $C$ accepts its input; this description is the
quadratic circuit construction's four-query PCPP description.
\end{definition}

\begin{theorem}[Positivity of the query count]
\lean{FABL.circuitQuadraticFourQuerySemanticPCPPReduction_queryCount_pos}
\label{FABL.circuitQuadraticFourQuerySemanticPCPPReduction_queryCount_pos}
\leanok
\uses{FABL.BooleanCircuit, FABL.BooleanCircuit.circuitQuadraticFourQueryPCPPDescription, FABL.circuitQuadraticFourQuerySemanticPCPPReduction}
\difficulty{1}
Let $n$ be a positive natural number and let $C$ be a Boolean circuit on $n$ input bits.
The query count in the four-query PCPP description assigned to $C$ by the semantic
reduction is positive.
\end{theorem}

\begin{definition}[Three-query circuit PCPP reduction]
\lean{FABL.quadraticCircuitSemanticPCPPReduction}
\statementsource{Analysis-of-Boolean-Functions-by-Ryan-ODonnell}{
  pdf-ffbe9c7b3306-p183-b003,
  pdf-ffbe9c7b3306-p198-b004
}
\label{FABL.quadraticCircuitSemanticPCPPReduction}
\leanok
\uses{FABL.SemanticPCPPReduction, FABL.circuitQuadraticFourQuerySemanticPCPPReduction, FABL.circuitQuadraticFourQuerySemanticPCPPReduction_queryCount_pos, FABL.semanticPCPPReductionToThreeQueries}
For every $n>0$ and every Boolean circuit $C$ on $n$ input bits, the reduction assigns a
PCPP description for the property that $C$ accepts its input, obtained by applying the
semantic conversion from four queries to three queries to the quadratic circuit's
four-query PCPP description.
\end{definition}

\begin{theorem}[Resources of the quadratic circuit reduction]
\lean{FABL.quadraticCircuitSemanticPCPPReduction_resources}
\proofsource{Analysis-of-Boolean-Functions-by-Ryan-ODonnell}{pdf-ffbe9c7b3306-p183-b003}
\label{FABL.quadraticCircuitSemanticPCPPReduction_resources}
\leanok
\uses{FABL.BooleanCircuit, FABL.BooleanCircuit.circuitQuadraticThreeQueryPCPPDescription_rejectionRate, FABL.quadraticCircuitSemanticPCPPReduction}
\difficulty{1}
For every positive integer $n$ and every Boolean circuit $C$ on $n$ input bits, the
semantic PCPP reduction associated with $C$ has query count $3$ and rejection rate
$1/12800000$.
\end{theorem}

\begin{theorem}[Proof length of the quadratic circuit reduction]
\lean{FABL.quadraticCircuitSemanticPCPPReduction_proofLength}
\label{FABL.quadraticCircuitSemanticPCPPReduction_proofLength}
\leanok
\uses{FABL.BooleanCircuit, FABL.BooleanCircuit.circuitQuadraticFourQueryPCPPSystem, FABL.BooleanCircuit.circuitQuadraticPCPPProofLength, FABL.BooleanCircuit.circuitQuadraticThreeQueryPCPPDescription_proofLength, FABL.PCPPSystem.descriptionSize, FABL.quadraticCircuitSemanticPCPPReduction}
\difficulty{1}
Let $n \ge 1$ and let $C$ be a Boolean circuit on $n$ input bits. The proof length of
the three-query semantic PCPP reduction associated with $C$ is $C$’s quadratic PCPP
proof length plus $64s$, where $s$ is the number of random seeds in the four-query PCPP
system for $C$.
\end{theorem}

\begin{theorem}[Proof length of the quadratic circuit reduction]
\lean{FABL.quadraticCircuitSemanticPCPPReduction_proofLength_le_explicit}
\label{FABL.quadraticCircuitSemanticPCPPReduction_proofLength_le_explicit}
\leanok
\uses{FABL.BooleanCircuit, FABL.BooleanCircuit.circuitQuadraticFourQueryPCPPSystem, FABL.BooleanCircuit.circuitQuadraticFourQueryPCPPSystem_descriptionSize_le, FABL.BooleanCircuit.circuitQuadraticPCPPProofLength, FABL.BooleanCircuit.circuitQuadraticPCPPProofLength_le_three_mul_two_pow_square, FABL.BooleanCircuit.circuitQuadraticVariableCount, FABL.PCPPSystem.descriptionSize, FABL.quadraticCircuitSemanticPCPPReduction, FABL.quadraticCircuitSemanticPCPPReduction_proofLength}
\difficulty{5}
Let $n$ be a positive integer and let $C$ be a Boolean circuit with $n$ input bits. The
proof length of the PCPP description assigned to $C$ by the quadratic circuit semantic
reduction is at most $2^{12(n+g+2)^2}$, where $g$ is the number of gates of $C$.
\end{theorem}

\begin{definition}[Quadratic circuit output-sensitive adapter]
\lean{FABL.QuadraticCircuitOutputSensitiveAdapter}
\label{FABL.QuadraticCircuitOutputSensitiveAdapter}
\leanok
\uses{FABL.BooleanCircuitEncoding, FABL.ThreeQueryOutputSensitiveAdapter, FABL.circuitQuadraticFourQuerySemanticPCPPReduction, FABL.circuitQuadraticFourQuerySemanticPCPPReduction_queryCount_pos}
Given a binary encoding of packed Boolean circuits, the quadratic circuit
output-sensitive adapter is the three-query output-sensitive adapter formed using that
encoding, the circuit's quadratic four-query semantic PCPP reduction, and the
reduction's positive query-count property.
\end{definition}

\begin{theorem}[Semantic output of the quadratic circuit adapter]
\lean{FABL.QuadraticCircuitOutputSensitiveAdapter.semantic_eq_canonical}
\label{FABL.QuadraticCircuitOutputSensitiveAdapter.semantic_eq_canonical}
\leanok
\uses{FABL.BooleanCircuit, FABL.BooleanCircuitEncoding, FABL.QuadraticCircuitOutputSensitiveAdapter, FABL.quadraticCircuitSemanticPCPPReduction}
\difficulty{1}
Let $C$ be a Boolean circuit with a positive number of input bits, and let an
output-sensitive adapter be built for a binary encoding of circuits. The semantic PCPP
description produced by the adapter for $C$ equals the canonical three-query PCPP
description obtained from the quadratic circuit construction.
\end{theorem}
